\chapter{模型选择与超参数}

\section{学习目标}

学完本章后，你应该能够：

\begin{itemize}
    \item 区分模型参数与超参数，并理解两者在训练流程中的角色不同；
    \item 解释为什么模型选择必须依赖验证集或交叉验证，而不能直接偷看测试集；
    \item 理解欠拟合、过拟合与模型复杂度之间的关系；
    \item 用一个明确的工作流完成候选模型比较、超参数选择和最终测试；
    \item 知道在数据量有限时，为什么交叉验证往往比单次切分更稳健。
\end{itemize}

\section{参数和超参数不是一回事}

在机器学习里，很多学生第一次遇到“模型选择”时，会把所有可调内容混在一起。实际上，至少要区分两类量：

\begin{itemize}
    \item \textbf{参数}：由训练过程根据数据学出来，例如线性回归的系数、神经网络的权重；
    \item \textbf{超参数}：由研究者在训练外部指定，例如多项式阶数、决策树深度、$k$ 近邻中的 $k$、正则化强度等。
\end{itemize}

参数回答的是“模型如何拟合这一批训练数据”；超参数回答的是“我们允许模型具有什么复杂度和归纳偏好”。

\section{为什么模型选择是一个独立问题}

如果你一口气尝试很多模型和很多超参数组合，那么你其实不只是“训练了一个模型”，而是在进行一轮\textbf{选择}。一旦有选择，就必须警惕：

\begin{itemize}
    \item 是根据什么标准做出的选择；
    \item 这个标准是否已经偷看了测试集；
    \item 你报告的最终分数是否被选择过程污染。
\end{itemize}

这就是为什么真实工作流里通常要把数据分成：

\begin{itemize}
    \item \textbf{训练集}：拟合模型参数；
    \item \textbf{验证集}：比较候选模型与超参数；
    \item \textbf{测试集}：在所有选择都冻结后，只做最后一次报告。
\end{itemize}

\section{复杂度太低会欠拟合，太高会过拟合}

超参数往往直接控制模型复杂度。例如：

\begin{itemize}
    \item 多项式回归的阶数越高，函数越灵活；
    \item 决策树的最大深度越大，划分边界越复杂；
    \item $k$ 近邻里的 $k$ 越小，模型越容易追随局部细节。
\end{itemize}

如果复杂度太低，模型连主要趋势都抓不住，会出现\textbf{欠拟合}；如果复杂度太高，模型会把噪声也当成规律，会出现\textbf{过拟合}。模型选择的任务，就是在这两者之间找到更合适的平衡。

\section{一个教学版恒星颜色--温度例子}

配套 notebook 使用 \nolinkurl{stellar_teff_model_selection_demo.csv}，它是一个小型的 Gaia 风格教学数据集，包含：

\begin{itemize}
    \item 颜色指标 \texttt{bp\_rp}；
    \item 参考有效温度 \texttt{teff\_k}。
\end{itemize}

这个例子有两个优点：

\begin{itemize}
    \item 颜色与温度之间确实存在明显但非线性的关系；
    \item 我们可以用\textbf{多项式阶数}作为超参数，直观地观察模型复杂度变化。
\end{itemize}

\section{单次验证集可以帮助你排除明显不合适的模型}

在 notebook 里，我们先把候选阶数设为：

\begin{examplebox}[候选超参数]
\begin{lstlisting}[style=pythonstyle]
degrees = [1, 2, 3, 5]
\end{lstlisting}
\end{examplebox}

随后先做一次固定的训练/验证切分，比较每个阶数的训练集 RMSE 与验证集 RMSE。你会看到：

\begin{itemize}
    \item 一次项模型偏简单，验证误差偏高；
    \item 五次多项式在训练集上误差很小，但验证误差明显恶化；
    \item 二次和三次模型都比一次项合理得多。
\end{itemize}

这一步已经足以建立一个重要直觉：\textbf{训练误差持续下降，不代表模型更好。}

\section{交叉验证是在给选择过程降噪}

但如果数据量不大，单次验证切分的结果可能对样本选择非常敏感。这个教学例子里就会出现一个很典型的现象：单次验证切分看起来可能更偏向三次模型，但 3 折交叉验证的平均 RMSE 会把二次模型推到更稳定的位置。

这正说明：

\begin{itemize}
    \item 单次切分更像一次局部观察；
    \item 交叉验证是在多个切分上重复比较，估计哪种超参数更稳；
    \item 当样本不大时，交叉验证往往能减少“偶然好运”的影响。
\end{itemize}

\section{测试集必须最后才出现}

完成超参数选择之后，我们再用训练集与验证集的合并数据重新拟合最终模型，并且只在这时看一次测试集结果。这一步的逻辑非常重要：

\begin{itemize}
    \item 训练集和验证集负责\textbf{选择}；
    \item 测试集负责\textbf{报告}。
\end{itemize}

如果你在测试集上反复比较不同阶数，再挑一个最好看的结果写进报告，那么测试集就已经不再是测试集，而是被你偷偷变成了验证集。

\section{网格搜索不是魔法，只是系统地比较候选}

很多库把模型选择封装成 \texttt{GridSearchCV} 或类似接口，但它做的事情本质上并不神秘：

\begin{enumerate}
    \item 列出候选超参数组合；
    \item 对每个组合做训练和验证；
    \item 根据选定指标排序；
    \item 选出最优配置。
\end{enumerate}

真正困难的不是会不会调用某个接口，而是：

\begin{itemize}
    \item 你为什么选择这个指标；
    \item 候选空间是否合理；
    \item 是否把预处理也包含在选择流程里；
    \item 是否避免了数据泄漏。
\end{itemize}

\section{一个值得记住的标准工作流}

\begin{enumerate}
    \item 先定义科学问题和评价指标；
    \item 划分训练、验证、测试，或在训练阶段使用交叉验证；
    \item 明确候选模型与超参数范围；
    \item 只用训练/验证信息完成选择；
    \item 冻结方案后，用训练加验证数据重新拟合；
    \item 最后只在测试集上评估一次。
\end{enumerate}

这套流程看起来慢一点，但它是在保护最终结论的可信度。

\section{搜索太多也会带来“看上去更好”的假象}

超参数搜索并不是越大越好。如果你在很小的数据集上试几十种模型、几百组超参数，最后总能挑出一个看起来特别漂亮的组合，但这未必说明你真的找到了更好的规律，也可能只是把噪声筛选成了“最佳模型”。

\begin{warnbox}[模型选择中的常见误区]
\begin{itemize}
    \item 用测试集挑超参数；
    \item 只看训练误差，不看验证误差；
    \item 一次切分结果很好，就认定模型已经稳定；
    \item 搜索范围极大，却没有解释为什么这些候选值得比较；
    \item 在预处理、特征选择和模型比较之间反复偷看测试集。
\end{itemize}
\end{warnbox}

\section{AI 助手如何帮助你}

在这一章，AI 助手很适合帮助你：

\begin{itemize}
    \item 把网格搜索、交叉验证和最终测试的职责分清楚；
    \item 检查你是否不小心在测试集上做了选择；
    \item 帮你整理候选超参数表格与实验记录。
\end{itemize}

但候选空间为何这样设、复杂度是否和科学问题匹配、最终该报告哪一版结果，仍然必须由你自己负责判断。

\section{本章小结}

模型选择不是训练结束后的附带步骤，而是决定结果是否可信的核心流程。超参数控制模型复杂度；验证集和交叉验证帮助你做选择；测试集只负责最后的独立检验。把这三件事分开，是从“会跑模型”走向“会做可靠科学推断”的关键一步。
